\section{Prinsip Aksesibilitas (a11y) dan Pengalaman Pengguna}

\textbf{Aksesibilitas web} memastikan bahwa aplikasi web dapat digunakan oleh orang dengan berbagai kemampuan, termasuk pengguna yang memakai pembaca layar, navigasi keyboard, atau perlu kontras warna yang cukup \cite{mdn-learn}. Sub-CPMK 1.3 menekankan penerapan prinsip aksesibilitas dan UX dalam desain antarmuka.

\textbf{WCAG} (Web Content Accessibility Guidelines) dari W3C WAI mendefinisikan prinsip: dapat dirasakan (\textit{perceivable}), dapat dioperasikan (\textit{operable}), dapat dipahami (\textit{understandable}), dan kokoh (\textit{robust}). Dalam implementasi HTML, penggunaan elemen semantik (\texttt{header}, \texttt{nav}, \texttt{main}, \texttt{form}, \texttt{label}), kontras warna yang memadai, dan alternatif teks untuk gambar adalah praktik dasar \cite{mdn-learn}.

Pengalaman pengguna (UX) yang baik mencakup konsistensi tampilan, umpan balik yang jelas (misalnya setelah submit form), dan pencegahan kesalahan (validasi, pesan error yang informatif). Menggabungkan aksesibilitas dengan prinsip UX menghasilkan produk yang inklusif dan mudah digunakan \cite{mdn-learn}.
